\documentclass[10pt,a4paper]{article}
\usepackage[margin=1.8cm,top=1.65cm,bottom=1.8cm]{geometry}
\usepackage[utf8]{inputenc}
\usepackage[T1]{fontenc}
\usepackage[spanish,es-tabla]{babel}
\usepackage{tgheros}
\renewcommand{\familydefault}{\sfdefault}
\usepackage{newtxmath}
\usepackage{anyfontsize}
\usepackage{microtype,xcolor,needspace,fancyhdr,lastpage}
\usepackage{hyperref}
\definecolor{accent}{HTML}{24566A}
\definecolor{muted}{HTML}{56636A}
\definecolor{rulecolor}{HTML}{A4BDC5}
\hypersetup{colorlinks=true,urlcolor=accent,linkcolor=accent,
  pdftitle={Camilo Martínez Estay - Curriculum vitae},pdfauthor={Camilo Martínez Estay}}
\setlength{\parindent}{0pt}
\setlength{\parskip}{0pt}
\setlength{\emergencystretch}{1em}
\linespread{1.06}
\pagestyle{fancy}
\fancyhf{}
\renewcommand{\headrulewidth}{0pt}
\fancyfoot[L]{\footnotesize\color{muted}Camilo Martínez Estay\quad /\quad Octubre de 2026}
\fancyfoot[R]{\footnotesize\color{muted}\thepage\,/\,\pageref*{LastPage}}
\setlength{\footskip}{24pt}
\raggedbottom

% Fechas a la izquierda; contenido en una sola columna de lectura.
\newlength{\datewidth}\setlength{\datewidth}{2.6cm}
\newlength{\colgap}\setlength{\colgap}{.45cm}
\newlength{\bodywidth}\setlength{\bodywidth}{\dimexpr\textwidth-\datewidth-\colgap\relax}
\newcommand{\cvsection}[1]{%
  \par\needspace{4\baselineskip}\vspace{3mm}%
  \hspace*{\dimexpr\datewidth+\colgap\relax}%
  {\color{accent}\bfseries\fontsize{10.5}{13}\selectfont\MakeUppercase{#1}}\par\vspace{2mm}}
\newcommand{\cvrow}[2]{%
  \par\noindent\begin{minipage}[t]{\datewidth}\vspace{0pt}
    \raggedright\small\color{muted}\strut #1
  \end{minipage}\hspace{\colgap}%
  \begin{minipage}[t]{\bodywidth}\vspace{0pt}\raggedright\strut #2\end{minipage}\par\vspace{1.5mm}}
\newcommand{\role}[1]{{\bfseries #1}}
\newcommand{\detail}[1]{\par\vspace{.5mm}#1}

\begin{document}
{\color{accent}\fontsize{25}{29}\selectfont\bfseries CAMILO MARTÍNEZ ESTAY}\par
\vspace{2mm}{\color{rulecolor}\hrule height 1pt}\vspace{2.5mm}
{\small Instituto de Matemáticas, Universidad de Valparaíso, Chile\par
\href{mailto:camilo.martinez@postgrado.uv.cl}{camilo.martinez@postgrado.uv.cl}\quad\textbar\quad
\href{https://cmartinezestay99.github.io}{cmartinezestay99.github.io}}

\cvsection{Intereses de investigación}
\cvrow{}{Álgebra, combinatoria algebraica, teoría de representaciones, grupos de Coxeter y álgebra computacional.}

\cvsection{Formación académica}
\cvrow{2025--presente}{\role{Magíster en Matemáticas}\par
Universidad de Valparaíso.
\detail{Director: Diego Lobos Maturana.}
\detail{Tesis: \emph{Invariantes combinatorios para clases de isomorfismo en la categoría $\mathcal{NT}$}. Finalización y defensa previstas: marzo--abril de 2027.}}
\cvrow{2019--2024}{\role{Licenciatura en Matemáticas}\par
Pontificia Universidad Católica de Valparaíso (PUCV).
\detail{Director: Felipe Riquelme.}
\detail{Tesina: \emph{La relación entre el grupo fundamental, los espacios de cubrimientos y las acciones de grupos}.}
\detail{Primer lugar de la generación (1/6). Promedio general: 6,3/7,0.}}

\cvsection{Experiencia de investigación}
\cvrow{2025--presente}{\role{Investigación de Magíster}\par
Universidad de Valparaíso.
\detail{Clasificación de álgebras graduadas en la categoría $\mathcal{NT}$ mediante invariantes combinatorios asociados a matrices triangulares inferiores. Estudio de particiones, composiciones, orden de dominancia, configuraciones combinatorias e invariantes matriciales para distinguir clases de isomorfismo.}
\detail{El proyecto se vincula con construcciones asociadas a grupos de Coxeter y teoría de representaciones.}}
\cvrow{2026--presente}{\role{Colaboración en álgebra computacional}\par
Con María Alejandra Álvarez, PUCV.
\detail{Estudio y clasificación computacional de álgebras Hom-Lie de dimensión cuatro. Reconstrucción y verificación en Python de cálculos simbólicos implementados originalmente en Maple; desarrollo de procedimientos reproducibles para comprobar condiciones algebraicas y casos de clasificación.}}
\cvrow{2020--2021}{\role{Ayudante de investigación}\par
Instituto de Matemáticas, PUCV.
\detail{Desarrollo de programas en Python y MATLAB para explorar ejemplos y conjeturas sobre álgebras graduadas en el contexto Blob--Soergel.}
\detail{Trabajo relacionado con el artículo de D. Lobos, \emph{Nil Graded Algebras Associated to Triangular Matrices and Their Applications to Soergel Calculus}, Journal of Pure and Applied Algebra \textbf{228} (2024), 107766. \href{https://doi.org/10.1016/j.jpaa.2024.107766}{doi:10.1016/j.jpaa.2024.107766}.}}

\cvsection{Manuscrito en preparación}
\cvrow{2026}{V. Delgado, D. Lobos, C. Martínez-Estay y R. Salama.
\emph{Combinatorial Invariants for Isomorphism Classes in the Category $\mathcal{NT}$}. En preparación.}

\newpage
\cvsection{Docencia}
\cvrow{Mar.--jul. 2025}{\role{Profesor · Universidad de Valparaíso}\par
Taller de Pensamiento Lógico Matemático II (LIMA215), carrera de Matemáticas.}
\cvrow{2022}{\role{Tutor académico · PUCV}\par
Tutorías de matemáticas de pregrado en programas universitarios de apoyo estudiantil.}
\cvrow{2019--2021}{\role{Ayudante de docencia · PUCV}\par
Cursos de pregrado de Cálculo, Álgebra, Geometría y Fundamentos de la Matemática.}

\cvsection{Becas y reconocimientos}
\cvrow{2026}{\role{Beca de Magíster Nacional ANID}, Agencia Nacional de Investigación y Desarrollo, Chile.}
\cvrow{2026}{\role{Apoyo financiero de la UNAM} para participar en la Segunda Escuela de Investigación en Geometría en Cuernavaca, México.}
\cvrow{2024}{\role{Reconocimiento al Mejor Egresado}, Licenciatura en Matemáticas, PUCV. Primer lugar entre seis egresados.}

\cvsection{Charlas, escuelas y congresos}
\cvrow{19 nov. 2026}{\role{Charla programada.} \emph{Invariantes para clases de isomorfismo en la categoría $\mathcal{NT}$}. XVI Jornadas de Matemática, Universidad de Valparaíso, Chile.}
\cvrow{2026}{Congreso Latinoamericano de Álgebra, Pontificia Universidad Católica del Perú, Lima, Perú.}
\cvrow{Jun.--jul. 2026}{Segunda Escuela de Investigación en Geometría en Cuernavaca, Universidad Nacional Autónoma de México, Cuernavaca, México.}
\cvrow{Oct. 2025}{V Escuela de Doctorado en Matemáticas, Valparaíso, Chile.}
\cvrow{2024--2025}{Escuelas de posgrado y encuentros académicos en las universidades de Concepción, La Frontera y Talca, Chile.}
\cvrow{2021--2025}{Participación en ediciones de los Encuentros de Álgebra y Teoría de Nudos, Chile.}
\cvrow{Abr. 2024}{Curso breve \emph{Introducción a las álgebras KLR}, dictado por Diego Lobos, Universidad de Valparaíso.}

\cvsection{Servicio académico y organización}
\cvrow{2026}{\role{Evaluador}, Campeonato Escolar de Matemáticas (CMAT), sede Aconcagua, Universidad de Valparaíso.}
\cvrow{Ago. 2025--presente}{\role{Coorganizador}, seminario Puerto de Matemática, Universidad de Valparaíso.}
\cvrow{2025}{\role{Comité de apoyo a la organización}, XV Jornadas de Matemática, Universidad de Valparaíso.}

\cvsection{Otra experiencia profesional}
\cvrow{2023--2024}{\role{Desarrollador de contenidos de Matemáticas · EDUGAMI S.A.}\par
Trabajo remoto. Diseño de problemas y soluciones detalladas para las pruebas de admisión universitaria PAES Matemática M1 y M2.}

\cvsection{Idiomas y herramientas}
\cvrow{}{Español: lengua materna. Inglés: preparación de certificación TOEFL/IELTS.
\detail{Python, MATLAB, \LaTeX{} y GeoGebra. Experiencia en reconstrucción en Python de procedimientos simbólicos desarrollados en Maple.}}
\end{document}
